\documentclass[aspectratio=169]{beamer}
\usetheme{VinUni}
\usepackage{vinuni-macros}
\usepackage{booktabs}
\usetikzlibrary{shapes.geometric}

%% ── Metadata ─────────────────────────────────────────────
\title[AICB · Ngày 17]{Memory Systems for Agents}
\subtitle{AICB-P2T3 · Ngày 17 · Chương 4 --- Agent Nâng Cao}
\author[Nguyễn Tiến Đồng]{Nguyễn Tiến Đồng}
\institute[VinUni]{VinUniversity}
\date[Tuần 4]{Phase 2 · Track 3 · Tuần 4}

%% ── Override titleframe: centered VinUniversity layout (day02 style) ──
\renewcommand{\titleframe}{%
  \begin{frame}[plain, noframenumbering]
    \begin{tikzpicture}[remember picture, overlay]
      %% Campus photo background (full bleed)
      \node[inner sep=0pt, anchor=center]
        at (current page.center)
        {\includegraphics[width=\paperwidth, height=\paperheight]{\vinunicampus}};
      %% Blue gradient overlay
      \fill[VinBlueDark, opacity=0.55]
        (current page.south west) rectangle (current page.north east);
      %% Bottom accent bar
      \fill[VinRed]
        (current page.south west) rectangle ([yshift=3pt]current page.south east);
      %% VinUniversity logo — centered at top
      \node[anchor=north, yshift=-18pt]
        at (current page.north)
        {\includegraphics[height=1.4cm]{\vinunilogowhite}};
      %% Title — centered, large
      \node[anchor=center, yshift=0.4cm,
            text=white,
            font=\fontsize{24}{30}\selectfont\bfseries,
            text width=0.80\paperwidth, align=center]
        (titleblock) at (current page.center)
        {\inserttitle};
      %% Subtitle — centered, italic
      \node[anchor=north, yshift=-10pt,
            text=white!75,
            font=\fontsize{11}{14}\selectfont\itshape,
            text width=0.75\paperwidth, align=center]
        (subtitleblock) at (titleblock.south)
        {\insertsubtitle};
      %% Red accent line
      \draw[VinRed, line width=2pt]
        ([yshift=-12pt]subtitleblock.south) ++(-2.5cm,0) -- ++(5cm, 0);
      %% Author — centered
      \node[anchor=north, yshift=-22pt,
            text=white,
            font=\fontsize{13}{17}\selectfont\bfseries]
        (authorblock) at (subtitleblock.south)
        {\insertauthor};
      %% Institute + date — centered at bottom
      \node[anchor=north, yshift=-8pt,
            text=white!70,
            font=\fontsize{9}{12}\selectfont]
        at (authorblock.south)
        {\insertinstitute\enspace·\enspace\insertdate};
    \end{tikzpicture}
  \end{frame}%
}

\begin{document}

%% ═══════════════════════════════════════════════════════════
%% TITLE + AGENDA
%% ═══════════════════════════════════════════════════════════

\titleframe
\hookframe{Tại sao agent của bạn quên mọi thứ sau mỗi conversation --- và làm sao fix nó đúng cách?}

%% ── Instructor Profile ─────────────────────────────────────
\begin{frame}{Giảng viên (Instructor)}
  \begin{columns}[T, onlytextwidth]
    \begin{column}{1.0\textwidth}
      \vspace{10pt}
      {\LARGE\bfseries\color{VinBlueDark} Nguyễn Tiến Đồng} \\[6pt]
      {\large\color{darkgray} Technical Director CMC AI}
      
      \vspace{16pt}
      \textbf{Kinh nghiệm chuyên môn:}
      \begin{itemize}
        \item 7+ năm nghiên cứu và phát triển các mô hình ML, AI.
        \item 5+ năm quản lý quy trình R\&D từ model tới AI Product thực tế.
      \end{itemize}
      
      \vspace{12pt}
      \textbf{Lĩnh vực ưu tiên (Focus Area):}
      \begin{itemize}
        \item AI product management
        \item Natural language processing (NLP)
        \item Computer vision
      \end{itemize}
    \end{column}
  \end{columns}
\end{frame}

\agendaframe

%% ═══════════════════════════════════════════════════════════
%% SECTION 1: TẠI SAO AGENT "QUÊN"?
%% ═══════════════════════════════════════════════════════════

\sectionframe{Tại sao Agent ``quên''?}
  {Context window có giới hạn --- và hầu hết agent không có bộ nhớ ngoài}

%% ── 1.1 Bài toán "quên" ──────────────────────────────────
\contentframe{Agent hiện tại --- Stateless by default}{
  \begin{columns}[c, onlytextwidth]
    \begin{column}{0.50\textwidth}
      \centering
      \begin{tikzpicture}[
        scale=0.8, transform shape,
        box/.style={rectangle, rounded corners=3pt, draw=VinBlue, fill=VinBlue!10,
          minimum width=2.2cm, minimum height=0.7cm, font=\small\bfseries, text=VinBlue},
        arrow/.style={-stealth, thick, VinBlue!70},
        xarrow/.style={-stealth, thick, VinRed!60, dashed},
      ]
        %% Session 1
        \node[box] (s1) at (0,0) {Session 1};
        \node[box] (llm1) at (3.5,0) {LLM};
        \draw[arrow] (s1) -- node[above, font=\scriptsize] {context} (llm1);

        %% Session 2
        \node[box] (s2) at (0,-1.5) {Session 2};
        \node[box] (llm2) at (3.5,-1.5) {LLM};
        \draw[arrow] (s2) -- node[above, font=\scriptsize] {context mới} (llm2);

        %% X between sessions
        \draw[xarrow] (llm1) -- node[right, font=\scriptsize, text=VinRed] {không truyền} (llm2);

        %% Label
        \node[font=\scriptsize\itshape, text=MidGray, below=0.4cm of llm2]
          {Mỗi session bắt đầu từ zero};
      \end{tikzpicture}
    \end{column}
    \begin{column}{0.46\textwidth}
      \begin{itemize}
        \item LLM \textbf{không có persistent state} --- mỗi API call là một request độc lập
        \item User nói ``tôi thích Python'' ở session 1 → session 2 agent \textbf{không nhớ}
        \item Conversation dài $>$50 turns → \textbf{hit context limit}
      \end{itemize}
      \vspace{4pt}
      \warningbox{Đây là vấn đề \#1 khi deploy agent thực tế: user kỳ vọng agent ``nhớ'' --- nhưng nó không.}
    \end{column}
  \end{columns}
}

%% ── 1.2 Analogy ─────────────────────────────────────────
\contentframe{Analogy: Bộ nhớ Agent giống não người}{
  \begin{columns}[c, onlytextwidth]
    \begin{column}{0.48\textwidth}
      \centering
      \begin{tikzpicture}[
        scale=0.65, transform shape,
        brain/.style={rectangle, rounded corners=5pt, minimum width=2.8cm, minimum height=0.8cm,
          font=\small\bfseries, text=white, align=center},
        arrow/.style={-stealth, thick, VinBlue!60},
      ]
        %% Human brain
        \node[font=\small\bfseries, text=VinBlue] at (0,2.2) {Não người};
        \node[brain, fill=VinRed!80] (wm) at (0,1.2) {Working Memory};
        \node[brain, fill=VinBlue] (lt) at (0,0) {Long-term Memory};
        \draw[arrow] (wm) -- node[right, font=\tiny, text=MidGray, text width=1.5cm, align=left] {consolidation} (lt);

        %% Agent
        \node[font=\small\bfseries, text=VinBlue] at (5,2.2) {Agent};
        \node[brain, fill=VinRed!80] (cw) at (5,1.2) {Context Window};
        \node[brain, fill=VinBlue] (ext) at (5,0) {External Store};
        \draw[arrow] (cw) -- node[right, font=\tiny, text=MidGray, text width=1.5cm, align=left] {persist facts} (ext);

        %% Equivalence
        \draw[dashed, MidGray, thick] (wm) -- node[above, font=\tiny, text=MidGray] {tương đương} (cw);
        \draw[dashed, MidGray, thick] (lt) -- node[above, font=\tiny, text=MidGray] {tương đương} (ext);
      \end{tikzpicture}
    \end{column}
    \begin{column}{0.48\textwidth}
      {\scriptsize
      \defbox{Context Window = RAM}{Nhanh, tạm thời, giới hạn dung lượng ($\sim$128K tokens)}
      \vspace{0pt}
      \defbox{External Store = Ổ cứng}{Chậm hơn, bền vững, gần như vô hạn (Redis, Vector DB)}
      \vspace{0pt}
      \vinbox{Insight}{Agent cần cả hai: \textbf{fast access} cho conversation hiện tại + \textbf{persistent storage} cho knowledge qua sessions}
      }
    \end{column}
  \end{columns}
}

%% ── 1.3 Thảo luận mở ──────────────────────────────────────
\contentframe{Thảo luận mở: Agent lý tưởng của bạn cần ``nhớ'' gì?}{
  \begin{columns}[T, onlytextwidth]
    \begin{column}{0.48\textwidth}
      \vinbox{Câu hỏi dành cho lớp}{
        Hãy nghĩ về dự án Agent của nhóm bạn. Nếu nó có trí nhớ vô hạn, nó nên nhớ những thông tin gì về Người dùng để tạo ra trải nghiệm tốt nhất?
      }
      \vspace{4pt}
      \warningbox{Ngược lại: Có những thông tin nào Agent \textbf{Tuyệt đối không nên} ghi nhớ để tránh rủi ro?}
    \end{column}
    \begin{column}{0.48\textwidth}
      \textbf{Hãy thử phân loại những gợi ý của bạn:}
      \begin{itemize}
        \item Đâu là thông tin \textbf{Bất biến} (Tên, Năm sinh, Nghề nghiệp...)?
        \item Đâu là thông tin \textbf{Động} (Trạng thái cảm xúc, Sự phàn nàn...)?
        \item Đâu là thông tin \textbf{Tri thức} (Coding style, Lỗi từng gặp...)?
      \end{itemize}
    \end{column}
  \end{columns}
}

%% ═══════════════════════════════════════════════════════════
%% SECTION 2: CONTEXT ENGINEERING FRAMEWORK
%% ═══════════════════════════════════════════════════════════

\sectionframe{Context Engineering Framework}
  {7 layers of context --- quản lý những gì agent ``thấy''}

%% ── 2.1 Seven Context Layers ──────────────────────────────
\contentframe{7 Context Layers --- Kiến trúc thông tin cho Agent}{
  \begin{center}
    \begin{tikzpicture}[
      scale=0.85, transform shape,
      layer/.style={rectangle, rounded corners=2pt, minimum height=0.55cm,
        font=\scriptsize\bfseries, text=white, align=center},
    ]
      %% Stacked layers (bottom to top)
      \node[layer, fill=DarkGray, minimum width=10cm] (L1) at (0,0)
        {Policy Context --- Guardrails, safety rules};
      \node[layer, fill=VinBlue!60, minimum width=9.2cm] (L2) at (0,0.7)
        {Tool Context --- Function outputs, API responses};
      \node[layer, fill=VinBlue!50, minimum width=8.4cm] (L3) at (0,1.4)
        {Retrieval Context --- RAG results, documents};
      \node[layer, fill=orange!70!black, minimum width=7.6cm] (L4) at (0,2.1)
        {Memory Context --- Recalled facts, episodes};
      \node[layer, fill=VinRed!70, minimum width=6.8cm] (L5) at (0,2.8)
        {User Context --- Preferences, history};
      \node[layer, fill=green!50!black, minimum width=6cm] (L6) at (0,3.5)
        {Task Context --- Objective, instructions};
      \node[layer, fill=VinBlue, minimum width=5.2cm] (L7) at (0,4.2)
        {System Context --- Persona, constraints};

      %% Arrow + label
      \draw[-stealth, thick, MidGray] (-5.5,0) -- (-5.5,4.2)
        node[midway, left, font=\scriptsize\itshape, text=MidGray, text width=1.2cm, align=center]
        {Priority khi trim};
    \end{tikzpicture}
  \end{center}
  \vspace{2pt}
  \begin{columns}[T, onlytextwidth]
    \begin{column}{0.48\textwidth}
      \vinbox{Quy tắc trim}{Khi gần token limit: \textbf{trim từ dưới lên}. Policy context trim cuối cùng (safety không bao giờ bỏ).}
    \end{column}
    \begin{column}{0.48\textwidth}
      \warningbox{Conflict resolution: user preference mâu thuẫn policy constraint → policy \textbf{luôn thắng}. Cần explicit rules trong system design.}
    \end{column}
  \end{columns}
}

%% ── 2.2 Token Budget ────────────────────────────────────
\statsframe{Token Budget --- Phân bổ context window}{
  \begin{columns}[c, onlytextwidth]
    \begin{column}{0.22\textwidth}\statcard{10\%}{Short-term\\memory}\end{column}
    \begin{column}{0.22\textwidth}\statcard{4\%}{Long-term\\facts}\end{column}
    \begin{column}{0.22\textwidth}\statcard{3\%}{Episodic\\memory}\end{column}
    \begin{column}{0.22\textwidth}\statcard{3\%}{Semantic\\knowledge}\end{column}
  \end{columns}
  \vspace{8pt}
  \bottomvinbox{Tổng: 20\% context window}{Phần còn lại dành cho system prompt, task instructions, tool outputs, và output generation. Vượt 20\% → context bị nhiễu, accuracy giảm.}
}

%% ── 2.3 Tình huống thảo luận ──────────────────────────────
\contentframe{Trò chơi Tình Huống: Quyết định Cắt Gọt Token}{
  \begin{columns}[T, onlytextwidth]
    \begin{column}{0.48\textwidth}
      \defbox{Bối cảnh}{
        Agent tư vấn của bạn có Context Limit nhỏ (để tối ưu chi phí). Bạn đang chạm mốc \textbf{95\% giới hạn Token}.
      }
      \vspace{2pt}
      \textbf{Trong Context hiện tại đang chứa:}
      \begin{enumerate}
        \item Policy: Luật cấm văng tục (15\%)
        \item Task: Định nghĩa luồng tư vấn (15\%)
        \item User Profile: Sở thích đã khai báo (10\%)
        \item Chat History: 30 phút chat qua lại (55\%)
      \end{enumerate}
    \end{column}
    \begin{column}{0.48\textwidth}
      \vinbox{Câu hỏi}{Khi có tin nhắn mới, hệ thống bị OOM (Out of Memory). Bạn phải Hi Sinh phần nào đầu tiên? Tại sao?}
      \vspace{4pt}
      \begin{itemize}
        \item \textit{A: Cắt bớt phần Luật an toàn (Policy)?}
        \item \textit{B: Xoá thông tin Sở thích?}
        \item \textit{C: Cắt đi các tin nhắn cũ nhất?}
        \item \textit{D: Gọi LLM thứ 2 để tóm tắt Chat History?}
      \end{itemize}
      \vspace{4pt}
      \textit{(Gợi ý: Nhắc lại hình chóp 7 Layers ưu tiên Trim)}
    \end{column}
  \end{columns}
}

%% ═══════════════════════════════════════════════════════════
%% SECTION 3: 4 LOẠI MEMORY
%% ═══════════════════════════════════════════════════════════

\sectionframe{Cognitive Memory Model --- 4 loại Memory}
  {Short-term, Long-term, Episodic, Semantic}

%% ── 3.1 Memory Taxonomy ──────────────────────────────────
\contentframe{4 loại Memory --- Cognitive Model cho AI Agents}{
  \begin{center}
    \begin{tikzpicture}[
      scale=0.78, transform shape,
      qbox/.style={rectangle, rounded corners=4pt, minimum width=4.8cm, minimum height=1.6cm,
        align=center, font=\scriptsize, text=DarkGray},
    ]
      %% 2x2 grid
      %% Top-left: Short-term
      \node[qbox, fill=VinRed!12, draw=VinRed!40] (st) at (-2.8,1.4)
        {{\small\bfseries\color{VinRed}Short-term (Working)}\\
         Context window buffer\\
         Nhanh, tạm thời, $\sim$128K tokens};

      %% Top-right: Long-term
      \node[qbox, fill=VinBlue!12, draw=VinBlue!40] (lt) at (2.8,1.4)
        {{\small\bfseries\color{VinBlue}Long-term (Declarative)}\\
         Redis, PostgreSQL\\
         User prefs, facts qua sessions};

      %% Bottom-left: Episodic
      \node[qbox, fill=orange!12, draw=orange!40] (ep) at (-2.8,-1.4)
        {{\small\bfseries\color{orange!70!black}Episodic}\\
         Log trải nghiệm có thứ tự\\
         ``Lần trước tôi đã làm gì?''};

      %% Bottom-right: Semantic
      \node[qbox, fill=green!12, draw=green!50!black] (sm) at (2.8,-1.4)
        {{\small\bfseries\color{green!50!black}Semantic}\\
         Embeddings + Vector DB\\
         Domain knowledge retrieval};

      %% Axis labels
      \node[font=\scriptsize\itshape, text=MidGray] at (-2.8,2.8) {Tạm thời};
      \node[font=\scriptsize\itshape, text=MidGray] at (2.8,2.8) {Bền vững};
      \node[font=\scriptsize\itshape, text=MidGray, rotate=90] at (-5.4,0) {Cá nhân};
      \node[font=\scriptsize\itshape, text=MidGray, rotate=90] at (5.4,0) {Tri thức};

      %% Divider lines
      \draw[MidGray!40, dashed] (0,-2.8) -- (0,2.8);
      \draw[MidGray!40, dashed] (-5.2,0) -- (5.2,0);
    \end{tikzpicture}
  \end{center}
  \vspace{-2pt}
  \bottomvinbox{Cognitive science mapping}{Working memory $\leftrightarrow$ Short-term \enspace|\enspace Declarative memory $\leftrightarrow$ Long-term \enspace|\enspace Episodic \& Semantic tương tự tên gọi}
}

%% ── 3.2 Short-term Memory Strategies ────────────────────
\contentframe{Short-term Memory --- Context Window Management}{
  \begin{columns}[c, onlytextwidth]
    \begin{column}{0.50\textwidth}
      \centering
      \begin{tikzpicture}[
        scale=0.78, transform shape,
        box/.style={rectangle, rounded corners=2pt, minimum width=2.6cm, minimum height=0.55cm,
          font=\scriptsize, text=DarkGray, draw=VinBlue!40, fill=VinBlue!8},
        arrow/.style={-stealth, thick, VinBlue!60},
        strat/.style={font=\scriptsize\bfseries, text=VinBlue, anchor=west},
      ]
        %% Strategy 1: Buffer
        \node[strat] at (-3.2,2.4) {Buffer};
        \node[box, minimum width=0.6cm, fill=VinBlue!15] at (-1.5,2.4) {M1};
        \node[box, minimum width=0.6cm, fill=VinBlue!15] at (-0.5,2.4) {M2};
        \node[box, minimum width=0.6cm, fill=VinBlue!15] at (0.5,2.4) {M3};
        \node[box, minimum width=0.6cm, fill=VinBlue!20] at (1.5,2.4) {M4};
        \node[box, minimum width=0.6cm, fill=VinRed!20, draw=VinRed!40] at (2.5,2.4) {M5};
        \node[font=\tiny\itshape, text=VinRed, anchor=west] at (3.1,2.4) {limit!};

        %% Strategy 2: Summary
        \node[strat] at (-3.2,1.2) {Summary};
        \node[box, minimum width=1.8cm, fill=orange!15, draw=orange!40] at (-0.7,1.2) {Summary};
        \node[box, minimum width=0.6cm, fill=VinBlue!15] at (1,1.2) {M4};
        \node[box, minimum width=0.6cm, fill=VinBlue!20] at (2,1.2) {M5};

        %% Strategy 3: Sliding Window
        \node[strat] at (-3.2,0) {Sliding};
        \node[box, minimum width=1.2cm, fill=green!15, draw=green!40] at (-1.1,0) {System};
        \node[box, minimum width=1cm, fill=orange!15, draw=orange!40] at (0.3,0) {Sum.};
        \node[box, minimum width=0.6cm, fill=VinBlue!15] at (1.3,0) {M4};
        \node[box, minimum width=0.6cm, fill=VinBlue!20] at (2.3,0) {M5};

        %% Best label
        \node[font=\tiny\bfseries, text=green!50!black, anchor=west] at (3.1,0) {Best!};
      \end{tikzpicture}
    \end{column}
    \begin{column}{0.46\textwidth}
      \textbf{3 strategies chính:}
      \begin{enumerate}
        \item \textbf{Buffer}: giữ tất cả --- đơn giản nhưng hit limit sau $\sim$50 turns
        \item \textbf{Summary}: LLM tóm tắt history cũ --- ổn định nhưng tốn thêm LLM calls
        \item \textbf{Sliding window}: system + summary + last K turns --- \textit{best tradeoff cho production}
      \end{enumerate}
      \vspace{4pt}
      \vinbox{Production tip}{Short-term memory nên chiếm tối đa \textbf{10\% context window}. Trim khi vượt --- keep recent N tokens, discard oldest.}
    \end{column}
  \end{columns}
}

%% ── 3.3 Long-term Memory với Redis ──────────────────────
\contentframe{Long-term Memory với Redis --- Persistent Cross-Session}{
  \begin{columns}[c, onlytextwidth]
    \begin{column}{0.48\textwidth}
      \centering
      \begin{tikzpicture}[
        scale=0.8, transform shape,
        box/.style={rectangle, rounded corners=3pt, draw=VinBlue, fill=VinBlue!10,
          minimum width=2cm, minimum height=0.6cm, font=\scriptsize\bfseries, text=VinBlue},
        db/.style={rectangle, rounded corners=1pt, draw=VinRed!60, fill=VinRed!8,
          minimum width=2cm, minimum height=0.6cm, font=\scriptsize\bfseries, text=VinRed},
        arrow/.style={-stealth, thick, VinBlue!60},
      ]
        %% Conversation flow
        \node[box] (conv) at (0,1.5) {Conversation};
        \node[box] (extract) at (0,0) {LLM Extract};
        \node[db] (redis) at (3.5,0) {Redis};
        \node[box] (next) at (3.5,1.5) {Next Session};

        \draw[arrow] (conv) -- node[left, font=\tiny, text=MidGray] {kết thúc} (extract);
        \draw[arrow] (extract) -- node[below, font=\tiny, text=MidGray] {key facts} (redis);
        \draw[arrow] (redis) -- node[right, font=\tiny, text=MidGray] {load profile} (next);

        %% TTL labels
        \node[font=\tiny\itshape, text=MidGray, text width=2.5cm, align=center, below=0.4cm of redis]
          {TTL: prefs 90d, facts 30d, sessions 7d};
      \end{tikzpicture}
    \end{column}
    \begin{column}{0.48\textwidth}
      \textbf{Ý tưởng:}
      \begin{itemize}
        \item Sau mỗi conversation, LLM \textbf{extract key facts} rồi store vào Redis với TTL
        \item Session mới: load user profile vào system prompt \textit{trước} khi user nói
      \end{itemize}
      \vspace{4pt}
      \vinbox{Redis data model}{
        \textbf{Hash}: preferences (language, style)\\
        \textbf{Set}: facts (``biết Python, học ML'')\\
        \textbf{List}: session history (recent)\\
        Tất cả \textbf{O(1) reads} --- production-ready
      }
    \end{column}
  \end{columns}
}

%% ── 3.4 Memory Management Flow ─────────────────────────
\contentframe{Memory Management Flow --- Buffer → Summarize → Store}{
  \begin{center}
    \begin{tikzpicture}[
      scale=0.85, transform shape,
      phase/.style={rectangle, rounded corners=4pt, minimum height=0.9cm,
        minimum width=2.2cm, align=center, font=\scriptsize\bfseries, text=white},
      arrow/.style={-stealth, thick, VinBlue!70},
      store/.style={rectangle, rounded corners=2pt, draw=VinBlue!50, fill=LightBlue!50,
        minimum width=1.8cm, minimum height=0.55cm, font=\scriptsize, text=VinBlue, align=center},
    ]
      %% Flow
      \node[phase, fill=VinRed!80] (buf) at (0,0) {1. Buffer\\{\tiny (Context Window)}};
      \node[phase, fill=orange!80!black] (sum) at (3.5,0) {2. Summarize\\{\tiny (LLM call)}};
      \node[phase, fill=VinBlue] (ext) at (7,0) {3. Extract\\{\tiny (Key facts)}};
      \node[phase, fill=green!50!black] (per) at (10.5,0) {4. Persist\\{\tiny (External store)}};

      \draw[arrow] (buf) -- (sum);
      \draw[arrow] (sum) -- (ext);
      \draw[arrow] (ext) -- (per);

      %% Stores below
      \node[store] (redis) at (9,-1.4) {Redis\\{\tiny long-term facts}};
      \node[store] (chroma) at (12,-1.4) {Chroma\\{\tiny semantic embeddings}};

      \draw[-stealth, dashed, green!50!black, thick] (per) -- (redis);
      \draw[-stealth, dashed, green!50!black, thick] (per) -- (chroma);

      %% Trigger
      \node[font=\tiny\itshape, text=MidGray, below=0.5cm of buf]
        {Trigger: token count $>$ threshold};
    \end{tikzpicture}
  \end{center}
  \vspace{4pt}
  \begin{columns}[T, onlytextwidth]
    \begin{column}{0.48\textwidth}
      \vinbox{Write-back rule}{Chỉ persist sau \textbf{task completion} --- không write giữa chừng để tránh inconsistent state}
    \end{column}
    \begin{column}{0.48\textwidth}
      \warningbox{Conflict resolution: long-term fact mâu thuẫn short-term info → \textbf{recency wins}, flag for review.}
    \end{column}
  \end{columns}
}

%% ═══════════════════════════════════════════════════════════
%% SECTION 4: IMPLEMENTATION DEEP-DIVE
%% ═══════════════════════════════════════════════════════════

\sectionframe{Implementation Deep-Dive}
  {Code-level: LangGraph nodes cho mỗi memory type}

%% ── 4.1 LangGraph Memory Architecture ──────────────────
\begin{frame}[fragile]{LangGraph Memory State --- Code-Level}
  \begin{columns}[T, onlytextwidth]
    \begin{column}{0.52\textwidth}
      \begin{codeframe}{python}
class MemoryState(TypedDict):
    messages: list[BaseMessage]
    user_profile: dict        # long-term
    episodes: list[dict]      # episodic
    semantic_hits: list[str]  # semantic
    memory_budget: int        # tokens left

# Memory router: chọn loại phù hợp
def retrieve_memory(state):
    query = state["messages"][-1].content
    return {
        "user_profile": redis.hgetall(uid),
        "episodes": find_similar(query, k=3),
        "semantic_hits": chroma.query(query),
    }
      \end{codeframe}
    \end{column}
    \begin{column}{0.44\textwidth}
      \textbf{Integration pattern:}
      \begin{enumerate}
        \item Node \texttt{load\_memory}: đọc 3 loại memory khi bắt đầu
        \item Inject vào system prompt theo priority
        \item Node \texttt{save\_memory}: ghi khi kết thúc
      \end{enumerate}
      \vspace{4pt}
      \vinbox{Priority order}{
        1. Short-term (gần nhất)\\
        2. Long-term facts (user prefs)\\
        3. Relevant episodes\\
        4. Semantic knowledge
      }
    \end{column}
  \end{columns}
\end{frame}

%% ── 4.2 Episodic Memory ────────────────────────────────
\contentframe{Episodic Memory --- Learning từ Past Trajectories}{
  \begin{columns}[c, onlytextwidth]
    \begin{column}{0.52\textwidth}
      \centering
      \begin{tikzpicture}[
        scale=0.8, transform shape,
        ep/.style={rectangle, rounded corners=3pt, draw=orange!60, fill=orange!8,
          minimum width=3.2cm, minimum height=0.55cm, font=\scriptsize, text=DarkGray, align=left},
        arrow/.style={-stealth, thick, orange!50},
      ]
        %% Episodes
        \node[ep] (e1) at (0,2) {Task: debug API};
        \node[ep] (e2) at (0,1) {Trajectory: tried X, Y};
        \node[ep] (e3) at (0,0) {Outcome: Y worked};
        \node[ep] (e4) at (0,-1) {Reflection: X fails vì...};

        \draw[arrow] (e1) -- (e2);
        \draw[arrow] (e2) -- (e3);
        \draw[arrow] (e3) -- (e4);

        %% New task
        \node[rectangle, rounded corners=3pt, draw=VinBlue, fill=VinBlue!10,
          minimum width=3cm, minimum height=0.55cm, font=\scriptsize\bfseries, text=VinBlue]
          (new) at (5.5,0.5) {New similar task};

        \draw[-stealth, dashed, VinBlue!60, thick] (e4.east) -- ++(0.5,0) |- (new.west)
          node[pos=0.25, below, font=\tiny\itshape, text=MidGray] {similarity search};
      \end{tikzpicture}
    \end{column}
    \begin{column}{0.44\textwidth}
      \textbf{Lưu tuple mỗi episode:}
      \begin{itemize}
        \item \texttt{(task, trajectory, outcome, reflection)}
        \item Agent biết: ``approach X đã fail vì Y trong task tương tự''
      \end{itemize}
      \vspace{4pt}
      \vinbox{Smart Forgetting}{
        \textbf{LRU}: xóa episode ít dùng nhất\\
        \textbf{Importance decay}: score giảm theo thời gian\\
        \textbf{Consolidation}: merge episodes tương tự\\
        Voyager-style: extract reusable strategy → skill library
      }
    \end{column}
  \end{columns}
}

%% ── 4.3 Semantic Memory ────────────────────────────────
\contentframe{Semantic Memory --- Vector DB cho Knowledge Retrieval}{
  \begin{columns}[c, onlytextwidth]
    \begin{column}{0.48\textwidth}
      \centering
      \begin{tikzpicture}[
        scale=0.78, transform shape,
        box/.style={rectangle, rounded corners=3pt, draw=green!50!black, fill=green!8,
          minimum width=2.2cm, minimum height=0.6cm, font=\scriptsize\bfseries,
          text=green!50!black},
        arrow/.style={-stealth, thick, green!40!black},
      ]
        \node[box] (doc) at (0,1.5) {Domain Docs};
        \node[box] (embed) at (0,0) {Embed};
        \node[box] (vdb) at (3.5,0) {Chroma DB};
        \node[box] (query) at (3.5,1.5) {Agent Query};
        \node[box] (result) at (6.5,0.75) {Top-K};

        \draw[arrow] (doc) -- (embed);
        \draw[arrow] (embed) -- node[below, font=\tiny, text=MidGray] {vectors} (vdb);
        \draw[arrow] (query) -- node[right, font=\tiny, text=MidGray] {cosine sim} (vdb);
        \draw[arrow] (vdb) -- (result);
      \end{tikzpicture}
    \end{column}
    \begin{column}{0.48\textwidth}
      \begin{itemize}
        \item Encode domain knowledge → embeddings → Chroma/Pinecone
        \item Query = task description → cosine similarity → top-k chunks
        \item Agent discover facts mới → add vào DB với metadata \texttt{(source, confidence, timestamp)}
      \end{itemize}
      \vspace{4pt}
      \vinbox{Self-improving}{Agent tự mở rộng knowledge base qua interactions --- incremental knowledge growth}
    \end{column}
  \end{columns}
}

%% ── 4.4 Unified Architecture ───────────────────────────
\contentframe{Memory Architecture --- Combining All 4 Types}{
  \begin{center}
    \begin{tikzpicture}[
      scale=0.82, transform shape,
      mem/.style={rectangle, rounded corners=3pt, minimum width=2cm, minimum height=0.6cm,
        font=\scriptsize\bfseries, text=white, align=center},
      arrow/.style={-stealth, thick, VinBlue!60},
    ]
      %% Agent in center
      \node[rectangle, rounded corners=5pt, draw=VinBlue, fill=VinBlue!10,
        minimum width=2.5cm, minimum height=1cm, font=\small\bfseries, text=VinBlue,
        align=center, text width=2.3cm]
        (agent) at (0,0) {Agent\\{\scriptsize retrieve(query)}};

      %% 4 memory types around it
      \node[mem, fill=VinRed!80] (st) at (-4,1.2) {Short-term};
      \node[mem, fill=VinBlue] (lt) at (4,1.2) {Long-term};
      \node[mem, fill=orange!80!black] (ep) at (-4,-1.2) {Episodic};
      \node[mem, fill=green!50!black] (sm) at (4,-1.2) {Semantic};

      %% Arrows
      \draw[arrow] (st) -- node[above, font=\tiny, text=MidGray] {priority 1} (agent);
      \draw[arrow] (lt) -- node[above, font=\tiny, text=MidGray] {priority 2} (agent);
      \draw[arrow] (ep) -- node[below, font=\tiny, text=MidGray] {priority 3} (agent);
      \draw[arrow] (sm) -- node[below, font=\tiny, text=MidGray] {priority 4} (agent);

      %% Merged output
      \node[rectangle, rounded corners=3pt, draw=VinBlue!50, fill=LightBlue!50,
        minimum width=2cm, minimum height=0.5cm, font=\scriptsize, text=VinBlue]
        (merged) at (0,-2.3) {Merged context → LLM};
      \draw[arrow] (agent) -- (merged);
    \end{tikzpicture}
  \end{center}
  \vspace{2pt}
  \warningbox{Unified interface: \texttt{retrieve(query, types=["all"])} trả về merged context từ cả 4 loại memory, đã trim theo token budget.}
}

%% ═══════════════════════════════════════════════════════════
%% SECTION 5: FRAMEWORKS & PRIVACY
%% ═══════════════════════════════════════════════════════════

\sectionframe{Frameworks chuyên dụng \& Privacy}
  {Mem0, Zep --- và khi nào dùng framework có sẵn}

%% ── 5.1 Mem0 & Zep ─────────────────────────────────────
\contentframe{Mem0 \& Zep --- Managed Memory Layers}{
  \begin{columns}[T, onlytextwidth]
    \begin{column}{0.48\textwidth}
      \vinbox{Mem0}{
        \begin{itemize}
          \item Auto-classify memory types
          \item Smart retrieval: relevance + recency ranking
          \item Claim: \textbf{90\% token reduction}, 91\% faster retrieval
          \item API-first, nhanh go-to-market
        \end{itemize}
      }
    \end{column}
    \begin{column}{0.48\textwidth}
      \vinbox{Zep}{
        \begin{itemize}
          \item Entity extraction + progressive summarization
          \item Tự build user knowledge graph qua sessions
          \item Multi-level summaries: turn → session → cross-session
          \item Giảm context size tối ưu
        \end{itemize}
      }
    \end{column}
  \end{columns}
  \vspace{2pt}
  \centering
  \renewcommand{\arraystretch}{1.1}
  {\footnotesize
  \rowcolors{2}{LightBlue!30}{white}
  \begin{tabular}{@{}l l l@{}}
    \toprule
    \textbf{Tiêu chí} & \textbf{Mem0 / Zep} & \textbf{Custom (Redis + Chroma)} \\
    \midrule
    Setup time   & Nhanh (API)      & Chậm (build from scratch) \\
    Control      & Hạn chế          & Full control \\
    Khi nào dùng & MVP, go-to-market & Production, đặc thù domain \\
    \bottomrule
  \end{tabular}
  }
}

%% ── 5.2 Privacy & GDPR ─────────────────────────────────
\contentframe{Quyền riêng tư, Bảo mật \& GDPR}{
  \begin{columns}[T, onlytextwidth]
    \begin{column}{0.48\textwidth}
      \defbox{Privacy-by-Design}{Mặc định \textbf{không lưu PII}. User phải explicit opt-in trước khi agent ghi nhớ thông tin cá nhân.}
      \vspace{6pt}
      \defbox{Right to be Forgotten}{User yêu cầu xóa → xóa \textbf{tất cả} memory entries liên quan → confirm deletion.}
    \end{column}
    \begin{column}{0.48\textwidth}
      \warningbox{Federated Forgetting: trong multi-agent system, deletion request phải \textbf{propagate} đến tất cả agents có copy.}
      \vspace{6pt}
      \vinbox{GDPR checklist}{
        \cmark Data minimization\\
        \cmark Purpose limitation\\
        \cmark Storage limitation (TTL)\\
        \cmark Consent management\\
        \cmark Deletion verification
      }
    \end{column}
  \end{columns}
}

%% ── 5.3 Tranh luận Đạo đức ─────────────────────────────────
\contentframe{Tranh luận: Dao hai lưỡi của ``Trí nhớ Agent''}{
  \begin{center}
    \textbf{``Ranh giới giữa sự Cá nhân hoá (Personalization) và Xâm phạm quyền riêng tư (Creepiness) nằm ở đâu?''}
  \end{center}
  \vspace{6pt}
  \begin{columns}[T, onlytextwidth]
    \begin{column}{0.48\textwidth}
      \defbox{Tình huống 1}{
        Một Agent trợ lý tiếng Anh tự động nhắc: \textit{``Hôm qua bạn bảo đang buồn vì vừa chia tay, hôm nay rảnh không, mình học từ vựng về chủ đề hàn gắn cảm xúc nhé?''}
      }
    \end{column}
    \begin{column}{0.48\textwidth}
      \defbox{Tình huống 2}{
        Agent công việc báo cáo lại: \textit{``Tôi nhận thấy bạn thường gõ sai phím và nổi cáu với AI lúc 11h đêm. Có vẻ bạn mệt rồi, tôi sẽ dời email gửi đi sang sáng mai.''}
      }
    \end{column}
  \end{columns}
  \vspace{8pt}
  \vinbox{Thảo luận}{Các hành vi trên là Thông Minh hay Đáng Sợ? Là một Kỹ sư AI, bạn sẽ thiết lập luồng \textbf{Xin phép (Opt-in)} cho tính năng lưu trữ kỉ niệm cá nhân này như thế nào?}
}

%% ═══════════════════════════════════════════════════════════
%% SECTION 6: DEMO & THỰC HÀNH
%% ═══════════════════════════════════════════════════════════

\sectionframe{Demo \& Thực hành}
  {Xem agent nhớ user preferences qua 3 sessions}

%% ── 6.1 Live Demo ──────────────────────────────────────
\demoframe{Agent nhớ User Preferences qua 3 Sessions}{
  \item \textbf{Session 1:} User nói ``tôi thích Python, không thích Java'' → agent ghi vào Redis
  \item \textbf{Session 2} (new process): Agent load memory → proactively suggest Python solution mà không cần hỏi lại
  \item \textbf{Session 3:} Agent recall episode ``user bị confused async/await'' → tự thêm explanation
  \item So sánh: agent có memory vs không memory --- response relevance, user satisfaction
}

%% ── 6.2 Lab ─────────────────────────────────────────────
\labframe{17}{Build Multi-Memory Agent với LangGraph}{
  Agent với full memory stack + benchmark report: so sánh agent có/không memory trên 10 multi-turn conversations
}{2 giờ}

\contentframe{Lab 17 --- Các bước thực hành}{
  \begin{enumerate}
    \item \textbf{Implement 4 memory backends}: ConversationBufferMemory (short-term), Redis (long-term), JSON episodic log, Chroma (semantic)
    \item \textbf{Build memory router}: chọn memory type phù hợp dựa trên query intent --- user preference vs factual recall vs experience recall
    \item \textbf{Context window management}: auto-trim khi gần limit, priority-based eviction theo 4-level hierarchy
    \item \textbf{Benchmark}: so sánh agent có/không memory trên 10 multi-turn conversations --- đo response relevance, context utilization, token efficiency
  \end{enumerate}
  \vspace{4pt}
  \vinbox{Deliverable}{GitHub repo + benchmark report: bảng so sánh metrics, memory hit rate analysis, token budget breakdown}
}

%% ═══════════════════════════════════════════════════════════
%% SUMMARY + CLOSING
%% ═══════════════════════════════════════════════════════════

\summaryframe{
  \takeaway{1}{Không có ``one size fits all'' --- production agent cần ít nhất \textbf{short-term + long-term}, thêm episodic/semantic tùy use case}
  \takeaway{2}{Memory retrieval quality quyết định agent quality --- bad retrieval = irrelevant context = wrong answer}
  \takeaway{3}{Memory write-back cần careful design: nhớ gì, khi nào ghi, xử lý conflict ra sao, TTL bao lâu}
  \takeaway{4}{Privacy không phải afterthought --- GDPR compliance cần thiết kế từ đầu (Privacy-by-Design)}
}

\previewframe{Ngày 18: Production RAG}
  {Agent đã có memory, tiếp theo cần knowledge retrieval tốt hơn --- tại sao RAG pipeline demo chạy tốt nhưng production chỉ đạt 60\%?}
  {\item Hoàn thành Lab 17: Multi-Memory Agent + benchmark
   \item Đọc: Anthropic ``Building Effective Agents'' (mục Context Engineering)}

\qaframe[Memory nào là ``must-have'' cho production agent? Khi nào thì dùng framework (Mem0, Zep) vs tự build?]

\closingframe{
  AICB-P2T3 · Ngày 17 · Memory Systems for Agents\\[4pt]
  \texttt{github.com/vinuni-aicb}\\[2pt]
  Liên hệ: \texttt{instructor@vinuni.edu.vn}
}

\end{document}
